\section{有限反冲如何把多边格缩成少态动力学}
\label{chap:feqo-bragg}

低反冲模型中每条边在作用时间内都几乎同相。反冲相位可分辨后，\cref{eq:feqo-exact-recoil-chain} 的 \(\delta_\ell\) 把它们分开；只有接近共振的通道可以持续积累振幅。设目标态为 \(|0\rangle,|1\rangle\)，其余态组成 \(Q\) 子空间，目标投影为 \(P\)。这是正交投影：\(P^2=P=P^\dagger\)，\(Q=I-P\)。

本节的两种近似各占一块区域。\cref{fig:recoil-regime-map} 在跃迁面积 \(\mathscr A_{\rm hop}\) 与反冲相位 \(\omega_{\rm r}T\) 张成的对数平面上比较它们的总变差误差：左图的参照是理想 Bessel 梯，\(\omega_{\rm r}T\) 小时整条横带都接近色标下限；右图的参照是共振二态极限，误差小的一角只出现在 \(\mathscr A_{\rm hop}\lesssim0.3\)、\(\omega_{\rm r}T\gtrsim1\) 处，\(\epsilon=0.05\) 的等高线先在 \(\mathscr A_{\rm hop}\approx0.3\) 附近近乎竖直，再随 \(\omega_{\rm r}T\) 增大向右弯折。两个极限占据相反的区域，中间地带既不是 Bessel 也不是二态；右图左上方明暗相间的水平条纹还说明二态误差随 \(\omega_{\rm r}T\) 并不单调下降。

\begin{figure}[htbp]
\centering
\includegraphics[width=0.94\linewidth]{recoil_regime_map.pdf}
\caption{有限反冲格点上的近似精度分区。横纵坐标分别为跃迁面积 \(\mathscr A_{\rm hop}\) 与反冲相位 \(\omega_{\rm r}T\)，色标为总变差误差的 \(\log_{10}\epsilon\)。左图以理想 Bessel 梯 \(J_\ell^2(2\mathscr A_{\rm hop})\) 为参照，右图以共振二态极限为参照；两种比较都在 \(0\leftrightarrow1\) 精确共振的调谐下进行，等高线取 \(\epsilon=0.05\)，仅用于区分近似精度区域。}
\label{fig:recoil-regime-map}
\end{figure}

对不随时间变的旋转绘景哈密顿量 \(H\)，本征方程 \(H|\psi\rangle=z|\psi\rangle\) 的 \(Q\) 分量给 \((z-QHQ)Q|\psi\rangle=QHP P|\psi\rangle\)。只要 \(z\) 不在 \(QHQ\) 的谱上，便能解出 \(Q|\psi\rangle=(z-QHQ)^{-1}QHP P|\psi\rangle\)，代回 \(P\) 分量：
\begin{equation}
H_{\rm eff}(z)=PHP+PHQ(z-QHQ)^{-1}QHP.
\label{eq:feqo-feshbach-effective}
\end{equation}
这是消元恒等式，不要求弱耦合。它也说明被删掉的通道并非没有作用：第二项会移动能级并改变有效耦合。

三态例子可把量级算清。令目标两态能量为 \(0,\hbar\Delta\)，第三态能量为 \(\hbar\Delta_Q\)，且三态 哈密顿量 除以 \(\hbar\) 后为
\[
\begin{pmatrix}
0&g&h_0\\g^*&\Delta&h_1\\h_0^*&h_1^*&\Delta_Q
\end{pmatrix}.
\]
当 \(|\Delta_Q|\gg |g|,|h_0|,|h_1|,|\Delta|\) 时，在 \(z\) 接近零的范围展开 \((z-\hbar\Delta_Q)^{-1}=-(\hbar\Delta_Q)^{-1}+O(z/(\hbar^2\Delta_Q^2))\)，得到
\begin{equation}
\frac{H_{\rm eff}}{\hbar}=
\begin{pmatrix}0&g\\g^*&\Delta\end{pmatrix}
-\frac1{\Delta_Q}
\begin{pmatrix}|h_0|^2&h_0h_1^*\\h_1h_0^*&|h_1|^2\end{pmatrix}
+O\!\left(\frac{\|h\|^2\max(|g|,|\Delta|)}{|\Delta_Q|^2}\right).
\label{eq:feqo-bragg-two-level}
\end{equation}
这里 \(\|h\|^2=|h_0|^2+|h_1|^2\)。被消去态的振幅约为 \(\|h\|/|\Delta_Q|\)，人口约为这个比值的平方。若该比值不小，两态方程连人口都不能准确预测。二阶对角项若对两个目标态不同，还会移动共振位置，所以一般不能只凭裸失谐 \(\Delta=0\) 判共振。

谱选择之所以能发生，是因为共振调谐本身把外部格点推远了。\cref{fig:recoil-two-level-selection} 画出把 \(0\leftrightarrow1\) 调成简并后的失谐尺度：\(\ell=0,1\) 简并（图中标注 resonant），而二阶色散下最近的两个外部格点 \(\ell=-1,2\) 各自失谐 \(\delta_{-1}\approx\delta_2\approx2\omega_{\rm r}\)，所以它们对目标双态只有虚贡献。图中只画出谱位置这一层机制，定量的截断判断仍要用完整 \(Q\) 子空间的预解式。

\begin{figure}[htbp]
\centering
\includegraphics[width=0.62\linewidth]{recoil_two_level_selection.pdf}
\caption{调谐 \(0\leftrightarrow1\) 共振后的谱选择。四条横线给出二阶色散下 \(\ell=-1,0,1,2\) 的失谐尺度：中间两态简并（resonant），两侧的 \(\ell=-1,2\) 各失谐约 \(2\omega_{\rm r}\)。图只说明谱选择机制。}
\label{fig:recoil-two-level-selection}
\end{figure}

现在把上一章的反冲算例真正代入，而不是只说“保留两态”。令 \(\omega_d=k_zv_0+\omega_r\)、\(\Delta_Q=2\omega_r\)，并先取所有相邻边的耦合为同一个实数 \(g\)。按 \(|-1\rangle,|0\rangle,|1\rangle,|2\rangle\) 排列，保留四态时的矩阵为
\[
\frac H\hbar=
\begin{pmatrix}
\Delta_Q&g&0&0\\
g&0&g&0\\
0&g&0&g\\
0&0&g&\Delta_Q
\end{pmatrix}.
\]
取 \(P\) 为中间两态、\(Q\) 为外侧两态，便有 \(QHQ/\hbar=\Delta_QI\)、\(PHQ/\hbar=gI\)。代入\cref{eq:feqo-feshbach-effective}，在 \(|z|\ll\hbar\Delta_Q\) 处展开，得到
\[
\frac{H_{\rm eff}}\hbar=
\begin{pmatrix}0&g\\g&0\end{pmatrix}
-\frac{g^2}{\Delta_Q}I
+O\!\left(\frac{g^3}{\Delta_Q^2}\right).
\]
此处两个目标态的二阶位移恰好相同，只产生共同相位，\emph{不}移动它们的相对共振；前面的三态模型中的差分位移不是所有几何都必然出现。若两侧耦合分别为 \(g_-\) 与 \(g_+\)，二阶对角位移改为 \(-|g_-|^2/\Delta_Q\) 与 \(-|g_+|^2/\Delta_Q\)，其差才移动共振。这个对称性核对也提醒我们：消元不仅要估计项的阶数，还要保留实际耦合结构。

约化后的两态问题给出的正是失谐 Rabi 振荡，进入振荡频率的是含外部格点自能的 \(\Delta_{\rm eff}\) 而不是裸失谐 \(\Delta\)。\cref{fig:rabi-detuning} 把三条不同失谐的曲线放在一起：\(\Delta=0\) 时振幅为 1、首个极大在 \(|\Omega|t=\pi\)；\(\Delta/|\Omega|=1\) 时广义 Rabi 频率 \(\sqrt{|\Omega|^2+\Delta^2}\) 变大，极大提前到 \(|\Omega|t=\pi/\sqrt2\)，峰值按 \(|\Omega|^2/(|\Omega|^2+\Delta^2)=1/2\) 下降；\(\Delta/|\Omega|=2\) 时峰值进一步压到 \(1/5\)，且振荡在 \(|\Omega|t/\pi\) 尺度上明显变密。这正是后面用脉冲面积而不是时间标定 Bragg 转移的原因。

\begin{figure}[htbp]
\centering
\includegraphics[width=0.76\linewidth]{rabi_detuning.pdf}
\caption{谱隔离后的失谐 Rabi 振荡。横轴为以 \(\pi\) 为单位的演化时间 \(|\Omega|t/\pi\)，纵轴为激发态布居 \(P_{\rm e}\)，三条曲线分别取 \(\Delta/|\Omega|=0,1,2\)。失谐使振荡变快、峰值按 \(|\Omega|^2/(|\Omega|^2+\Delta^2)\) 下降。}
\label{fig:rabi-detuning}
\end{figure}

四态矩阵是为了让消元可手算，并非声称原无限梯只有四个态。若所有更远边也保留，它们要从目标态至少经过两次相邻跳跃才到达，所以在最近失谐有下界且 \(|g|/\omega_r\ll1\) 时，不改变刚才的二阶结果；接近其他高阶共振或耦合不小，则必须返回\cref{eq:feqo-exact-recoil-chain} 数值求解。

这里还要区分两类误差：截断只是不给全部格点，而格点自身的曲率会让无限梯上的整条分布一起变形。\cref{fig:pinem-recoil-breaking} 把曲率相位 \(\chi=\omega_{\rm curv}T\) 从零打开，取 \(\beta=2.5\)：\(\chi=0\) 时分布是理想 \(J_\ell^2(2|\beta|)=J_\ell^2(5)\)，在 \(\ell=\pm4\) 处约 0.153、\(\ell=\pm3\) 处约 0.133，而 \(\ell=\pm2\) 处几乎为零（约 0.002）；\(\chi=0.20\) 时 \(\ell=\pm4\) 与 \(\ell=\pm3\) 被压成约 0.148 的平肩，\(\ell=\pm2\) 的极小开始被填到 0.007；\(\chi=0.40\) 时最大值移到 \(\ell=\pm3\)（约 0.184，反而超过理想值），\(\ell=\pm4\) 降到 0.113，\(\ell=\pm5\) 从 0.068 塌到 0.020。峰位与峰高的这种重新分配说明，一旦曲率不可忽略，理想 Bessel 梯就不是“再多加几阶”能修正的对象。

\begin{figure}[htbp]
\centering
\includegraphics[width=0.76\linewidth]{pinem_recoil_breaking.pdf}
\caption{曲率相位 \(\chi=\omega_{\rm curv}T\) 对边带分布的影响（\(\beta=2.5\)）。实线为理想 Bessel 梯 \(J_\ell^2(2|\beta|)\)，另外两条曲线是加入二次曲率项 \(\chi\ell^2\) 后的精确幺正演化。\(\chi\) 增大时峰位内移、\(\ell=\pm2\) 处的极小被填高，理想梯的峰高比例不再成立。}
\label{fig:pinem-recoil-breaking}
\end{figure}

纵向光场交换以 \(p_z\to p_z+\hbar k_z\) 为主，反冲尺度为 \(\omega_r=\hbar k_z^2/(2\gamma_0^3m)\)。横向驻波若沿 \(x\) 周期，连接的是 \(p_x\to p_x+\hbar G\)，中心电子仍沿 \(z\) 前进；在 \(p_x\simeq0\) 处横向色散的二次系数是 \(1/(\gamma_0m)\)，对应 \(\omega_{r,\perp}=\hbar G^2/(2\gamma_0m)\)。二者的有效质量相差 \(\gamma_0^2\)，不可共用一个未说明方向的反冲频率。

若 \(N^2\omega_rT\ll1\)，实际占据的多条边不可分辨，出现薄光栅或 Raman--Nath 极限；若一条边满足 \(|\delta_{\ell+1}-\delta_\ell|T/\hbar\lesssim1\)，而所有其余边满足 \(|g|/|\Delta_{\rm gap}|\ll1\)，则\cref{eq:feqo-bragg-two-level} 描述 Bragg 少态区。两组条件都不满足时应数值求解原反冲格。

“耦合不小会怎样”可以直接测出来。取相邻格点耦合 \(g=\hbar\Omega/2\)，只保留两态的预言是 \(\sin^2(|\Omega|T/2)\)；\cref{fig:electron-bragg-convergence} 用完整反冲格的含时解（TDSE）检验它。左图扫过矩形脉冲的脉冲面积：\(r_\Omega=|\Omega|/\omega_{\rm r}=0.2\) 时曲线与二态虚线几乎重合，名义 \(\pi\) 脉冲处 \(P_1\approx0.995\)；\(r_\Omega=1\) 时峰值降到约 0.92 且峰位向小面积一侧移动；\(r_\Omega=5\) 时 \(P_1\) 最高只有约 0.35，要等到 \(|\Omega|T=2\pi\) 附近才回升到 0.75。右图把脉冲固定在名义 \(\pi\) 处扫描驱动比：\(r_\Omega\lesssim0.5\) 时 \(P_1\) 保持在 1，泄漏 \(P_{\rm leak}\) 从 \(r_\Omega\approx0.5\) 起明显上升，到 \(r_\Omega=7\) 时 \(P_1\) 只剩约 0.1 而 \(P_{\rm leak}\) 达到约 0.8。所以“只保留两个通道”与前面的 \(g/\omega_{\rm r}\ll1\) 是同一个条件。

\begin{figure}[htbp]
\centering
\includegraphics[width=0.94\linewidth]{electron_bragg_convergence.pdf}
\caption{矩形脉冲下反冲格转移的 TDSE 收敛性检验。左图取不同驱动比 \(r_\Omega=|\Omega|/\omega_{\rm r}\)，比较 \(0\to1\) 转移概率 \(P_1\) 随脉冲面积 \(|\Omega|T/\pi\) 的变化，虚线为零泄漏的两态预言 \(\sin^2(|\Omega|T/2)\)；右图在名义 \(\pi\) 脉冲处给出目标概率 \(P_1\) 与泄漏概率 \(P_{\rm leak}\) 随 \(r_\Omega\) 的变化。计算保留双侧格点 \(|\ell|\le25\)。}
\label{fig:electron-bragg-convergence}
\end{figure}

最后回到贯穿全节的二阶色散。\cref{fig:recoil-exact-vs-quadratic} 用 20 keV、800 nm 的参数，把相邻跃迁频率的精确相对论值 \([\mathcal E(p+(\ell+1)\hbar\kappa)-\mathcal E(p+\ell\hbar\kappa)]/\hbar\) 与二阶反冲近似 \(v\kappa+(2\ell+1)\hbar\kappa^2/(2\gamma_0^3m)\) 各自减去公共平移 \(v\kappa\) 后画在一起。在 \(|\ell|\le12\) 内两条曲线不可区分：\(\ell=-12\) 与 \(\ell=+12\) 处分别约为 \(-158\) 与 \(+171\) GHz，最大偏差不到 0.01 GHz，相对量级约 \(5\times10^{-5}\)。这一参数下 \(\hbar\kappa/p_0\approx4\times10^{-5}\)，所以二阶色散确实够用；该图同时也提醒，“二阶够用”是数值结论而非恒等式，\(\hbar\kappa/p_0\) 变大时必须重算。

\begin{figure}[htbp]
\centering
\includegraphics[width=0.76\linewidth]{recoil_exact_vs_quadratic.pdf}
\caption{相邻跃迁频率的精确相对论值与二阶反冲近似（20 keV 电子、800 nm 光）。横轴为格点编号 \(\ell\)，纵轴为扣去公共平移 \(v\kappa\) 后的频率差 \(\delta f_\ell\)（GHz）。两条曲线在 \(|\ell|\le12\) 内不可区分，最大偏差约 0.01 GHz。}
\label{fig:recoil-exact-vs-quadratic}
\end{figure}

\begin{exercise}
令 \(h_0=h_1=h\)、\(g=0\)，求消去第三态后产生的目标态间耦合及两个能级的共同移动，并检验 \(h/\Delta_Q\to0\) 极限。
\end{exercise}
